\section*{Capítulo \ref{ch:g10:signals-and-waves} --- Señales y ondas}

\begin{solution}{exo:g10:signals-and-waves:1}
$f = 1/T$: (a) $1/\qty{0.020}{s} = \qty{50}{Hz}$;
(b) $1/\qty{4.0e-3}{s} = \qty{250}{Hz}$;
(c) $1/\qty{5.0e-5}{s} = \qty{20}{kHz}$.
\end{solution}

\begin{solution}{exo:g10:signals-and-waves:2}
$T = 1/f$: (a) \qty{20}{ms}; (b) $1/440 = \qty{2.3}{ms}$;
(c) $1/\qty{2.4e9}{Hz} = \qty{0.42}{ns}$.
\end{solution}

\begin{solution}{exo:g10:signals-and-waves:3}
\qty{12}{Hz}: infrasonido. \qty{440}{Hz} y \qty{18}{kHz}: audibles (el
segundo, solo para oídos jóvenes). \qty{40}{kHz} y \qty{5}{MHz}:
\omterm{def:g10:signals-and-waves:ultrasound}{ultrasonidos}.
\end{solution}

\begin{solution}{exo:g10:signals-and-waves:4}
$T = 5.0 \times \qty{2}{ms} = \qty{10}{ms}$, luego $f = \qty{100}{Hz}$;
\omterm{def:g10:signals-and-waves:amplitude}{amplitud} $4.0 \times \qty{0.5}{V} = \qty{2.0}{V}$.
\end{solution}

\begin{solution}{exo:g10:signals-and-waves:5}
$d = 340 \times 6.0 = \qty{2040}{m} \approx \qty{2.0}{km}$. La luz
recorre esa distancia en
$2040/(3.00\times10^{8}) \approx \qty{7}{\micro s}$, un millón de veces
menos que \qty{6.0}{s}: el destello marca el instante del impacto.
\end{solution}

\begin{solution}{exo:g10:signals-and-waves:6}
El \omterm{def:g10:signals-and-waves:sound}{sonido} baja \emph{y vuelve}, recorriendo $2d$:
$d = \frac{1500 \times 0.60}{2} = \qty{450}{m}$. Sobre la fosa:
$t = 2d/v = 2 \times 1200/1500 = \qty{1.6}{s}$.
\end{solution}

\begin{solution}{exo:g10:signals-and-waves:7}
$d_1 = \frac{3.00\times10^{8} \times 2.4\times10^{-4}}{2} =
\qty{36.0}{km}$ y $d_2 = \qty{35.7}{km}$: el avión recorrió
\qty{300}{m} en \qty{1.0}{s}, luego
$v = \qty{300}{m/s} \approx \qty{1100}{km/h}$, y se acerca.
\end{solution}

\begin{solution}{exo:g10:signals-and-waves:8}
$T = \qty{20}{mm}/\qty{25}{mm/s} = \qty{0.80}{s}$, luego
$f = \qty{1.25}{Hz}$: $75$ latidos por minuto. Para $100$ latidos por
minuto, $T = \qty{0.60}{s}$, es decir
$25 \times 0.60 = \qty{15}{mm}$: la taquicardia aparece por debajo de
\qty{15}{mm}.
\end{solution}

\begin{solution}{exo:g10:signals-and-waves:9}
$T = 1/200 = \qty{5.0}{ms}$; dos ciclos duran \qty{10}{ms}, repartidos
en $10$ divisiones: \omterm{def:g10:signals-and-waves:oscilloscope}{base de tiempos} de \qty{1}{ms} por división. La
cresta tiene que caber en $4$ divisiones: \omterm{def:g10:signals-and-waves:oscilloscope}{ganancia} de al menos
$3.0/4 = \qty{0.75}{V}$ por división --- se toma \qty{1}{V} por división
(cresta a $3.0$ divisiones).
\end{solution}

\begin{solution}{exo:g10:signals-and-waves:10}
$d = v\,t/2$: pared anterior,
$\frac{1540 \times 4.0\times10^{-5}}{2} = \qty{3.1}{cm}$; pared
posterior, $\frac{1540 \times 6.0\times10^{-5}}{2} = \qty{4.6}{cm}$;
grosor, unos \qty{1.5}{cm}.
\end{solution}

\begin{solution}{exo:g10:signals-and-waves:11}
\omterm{def:g10:signals-and-waves:sound}{Sonido}: $30/340 = \qty{0.088}{s}$. Radio:
$3.0\times10^{6}/(3.00\times10^{8}) = \qty{0.010}{s}$. El oyente de
radio, que está lejísimos, oye antes el tambor, por unos \qty{78}{ms}.
\end{solution}

\begin{solution}{exo:g10:signals-and-waves:12}
Mientras emite durante \qty{3.0}{ms}, el murciélago está sordo a los
\omterm{rem:g10:signals-and-waves:honest}{ecos}, es decir, a todo lo que esté a menos de
$d = \frac{340 \times 3.0\times10^{-3}}{2} = \qty{0.51}{m}$. Polilla a
\qty{1.7}{m}: $t = 2 \times 1.7/340 = \qty{10}{ms}$. Al acercarse, el
\omterm{rem:g10:signals-and-waves:honest}{eco} vuelve cada vez antes; el pulso tiene que terminar antes de que
llegue, así que hay que acortarlo.
\end{solution}

\begin{solution}{exo:g10:signals-and-waves:13}
$d = \frac{340 \times 1.5}{2} = \qty{255}{m}$. A \qty{155}{m}:
$t = 2 \times 155/340 = \qty{0.91}{s}$. El \omterm{rem:g10:signals-and-waves:honest}{eco} y la palmada se funden
cuando $t < \qty{0.1}{s}$, es decir a menos de
$d < \frac{340 \times 0.1}{2} = \qty{17}{m}$.
\end{solution}

\begin{solution}{exo:g10:signals-and-waves:14}
\emph{1.} $T = 3.0 \times \qty{5}{ms} = \qty{15}{ms}$,
$f = 1/0.015 \approx \qty{67}{Hz}$.
\emph{2.} $\qty{15}{ms}/\qty{2}{ms} = 7.5$ divisiones.
\emph{3.} La pantalla muestra \qty{50}{ms} y después \qty{20}{ms}: $3$
ciclos completos ($50/15 = 3.3$) y luego $1$ ($20/15 = 1.3$).
\end{solution}

\begin{solution}{exo:g10:signals-and-waves:15}
\emph{1.} $t = 2 \times 1.2\times10^{6}/(3.00\times10^{8}) =
\qty{8.0}{ms}$.
\emph{2.} $2 \times 1.2\times10^{6}/(2.0\times10^{8}) = \qty{12}{ms}$.
\emph{3.} La fibra no va en línea recta entre las dos máquinas, y cada
enrutador y cada servidor del camino añaden retraso de procesamiento.
\end{solution}

\begin{solution}{pb:g10:signals-and-waves:1}
\textbf{1.} $d = 340 \times 9.0 = \qty{3060}{m} \approx \qty{3.1}{km}$.

\textbf{2.} $3060/(3.00\times10^{8}) \approx \qty{10}{\micro s}$, un
millón de veces menos que \qty{9.0}{s}: el destello marca el instante
del impacto.

\textbf{3.} En \qty{3}{s} el \omterm{def:g10:signals-and-waves:sound}{sonido} recorre
$340 \times 3 = \qty{1020}{m} \approx \qty{1}{km}$: los segundos
divididos entre tres dan kilómetros.

\textbf{4.} $d = 340 \times 6.0 = \qty{2040}{m}$. La tormenta se ha
acercado \qty{1020}{m} en \qty{300}{s}:
$v = \qty{3.4}{m/s} \approx \qty{12}{km/h}$, rumbo al barco.

\textbf{5.} El extremo cercano se oye al cabo de
$2000/340 = \qty{5.9}{s}$ y el lejano al cabo de
$5000/340 = \qty{14.7}{s}$: el retumbo dura unos \qty{8.8}{s}.

\textbf{6.} $d = 1500 \times 0.90/2 = \qty{675}{m}$.

\textbf{7.} $d = 1500 \times 0.20/2 = \qty{150}{m}$.

\textbf{8.} Dos obstáculos: un banco de peces a
$1500 \times 0.16/2 = \qty{120}{m}$ y el fondo a \qty{150}{m}. El banco
nada \qty{30}{m} por encima del fondo.

\textbf{9.} El \omterm{rem:g10:signals-and-waves:honest}{eco} tiene que volver antes del pulso siguiente:
$d_{\max} = 1500 \times 0.50/2 = \qty{375}{m}$. Desde aguas más
profundas, el \omterm{rem:g10:signals-and-waves:honest}{eco} llega \emph{después} del siguiente pulso y se le
atribuye a él: la pantalla muestra un fondo falso, mucho menos profundo
de lo real.

\textbf{10.} En el aire, $340 \times 0.60/2 = \qty{102}{m}$ en lugar de
\qty{450}{m}. Un retraso se convierte en distancia solo cuando se sabe
cuál es el portador y su velocidad en el medio atravesado.

\textbf{11.} $f = 1/0.86 = \qty{1.16}{Hz}$, es decir
$1.16 \times 60 \approx 70$ latidos por minuto.

\textbf{12.} $25 \times 0.86 = \qty{21.5}{mm}$.

\textbf{13.} $T = 15/25 = \qty{0.60}{s}$: $100$ latidos por minuto.

\textbf{14.} $T = 60/50 = \qty{1.2}{s}$; separación
$25 \times 1.2 = \qty{30}{mm}$.

\textbf{15.} No: el \omterm{def:g10:signals-and-waves:period}{periodo} se desplaza de un latido a otro con el
esfuerzo, la respiración y el estrés. El médico lee el \omterm{def:g10:signals-and-waves:period}{periodo} (la
\omterm{def:g10:signals-and-waves:frequency}{frecuencia}), su regularidad y la forma de cada ciclo --- el diagnóstico
está en la \omterm{def:g10:signals-and-waves:signal}{señal}.

\textbf{16.} Radio: $5.0\times10^{4}/(3.00\times10^{8}) =
\qty{0.17}{ms}$ --- instantáneo para los sentidos humanos. \omterm{def:g10:signals-and-waves:sound}{Sonido}:
$50\,000/340 \approx \qty{147}{s}$, dos minutos y medio.

\textbf{17.} $t = 2.02\times10^{7}/(3.00\times10^{8}) = \qty{67}{ms}$.

\textbf{18.} $c \times \qty{1.0}{\micro s} = \qty{300}{m}$. Para una
posición con \qty{3.0}{m} de error:
$3.0/(3.00\times10^{8}) = \qty{10}{ns}$ --- y por eso los \omterm{def:g10:universal-gravitation:satellite}{satélites} GPS
llevan relojes atómicos.

\textbf{19.} Los \omterm{def:g10:universal-gravitation:satellite}{satélites} están en el vacío y el \omterm{def:g10:signals-and-waves:sound}{sonido} necesita un
medio: no saldría ninguna \omterm{def:g10:signals-and-waves:signal}{señal} del \omterm{def:g10:universal-gravitation:satellite}{satélite}. (Y aun concediendo aire en
todo el trayecto, \qty{20200}{km} a \qty{340}{m/s} son más de
\qty{16}{h} --- una posición con horas de retraso.)

\textbf{20.} \qty{340}{m/s} (\omterm{def:g10:signals-and-waves:sound}{sonido} en el aire), \qty{1500}{m/s} (\omterm{def:g10:signals-and-waves:sound}{sonido}
en agua de mar) y $c = \qty{3.00e8}{m/s}$ (radio y luz); la ley es
$d = v\,t$, dividida entre dos para un \omterm{rem:g10:signals-and-waves:honest}{eco}. La costumbre: saber cuál es
el portador, saber su velocidad en el medio y cronometrar el retraso ---
un retraso es una distancia disfrazada.
\end{solution}
